\section{Detailed Derivation of the ELBO}\label{sec:elbo_derivation}

\begin{story}[Notation]
As in the main text, $p_{\theta}(\mathbf{z}_{1} \mid \mathbf{z}_{0},
\star)$ is the prior on $\mathbf{z}_1$, and $q_{\phi}(\mathbf{z}_{1}
\mid \mathbf{z}_{0}, \star)$ is the first factor of the inference
network, with $\mathbf{z}_0$ fixed or learned. To lighten the formulas,
$q$ stands for $q_\phi(\cdot \mid \mathbf{x}_{1:T}, \mathbf{u}_{1:T})$
and its factors, and $p$ for $p_\theta$.
\end{story}

\subsection*{Step 1: the bound}

As for the VAE (Chapter~\ref{chap:vae}), Jensen's inequality gives
\begin{equation*}
    \log p(\mathbf{x}_{1:T} \mid \mathbf{u}_{1:T}) \;\geq\;
    \mathcal{L} := \E_{q(\mathbf{z}_{1:T})}\left[\log
    \frac{p(\mathbf{x}_{1:T}, \mathbf{z}_{1:T} \mid \mathbf{u}_{1:T})}
    {q(\mathbf{z}_{1:T})}\right],
\end{equation*}
with a gap equal to $\KL\big(q(\mathbf{z}_{1:T}) \,\|\,
p(\mathbf{z}_{1:T} \mid \mathbf{x}_{1:T}, \mathbf{u}_{1:T})\big)$.

\subsection*{Step 2: substitute the factorisations}

With the generative model \eqref{eq:generative_model} and the inference
model \eqref{q_phi},
\begin{equation*}
    \log \frac{p(\mathbf{x}_{1:T}, \mathbf{z}_{1:T} \mid \mathbf{u}_{1:T})}
    {q(\mathbf{z}_{1:T})} =
    \sum_{t=1}^T \log p(\mathbf{x}_t \mid \mathbf{z}_t)
    \;-\; \sum_{t=1}^T \underbrace{\log \frac{q(\mathbf{z}_t \mid
    \mathbf{z}_{t-1}, \mathbf{x}_{t:T}, \mathbf{u}_{t:T})}
    {p(\mathbf{z}_t \mid \mathbf{z}_{t-1}, \mathbf{u}_t)}}_{=:\ r_t(
    \mathbf{z}_{t-1}, \mathbf{z}_t)}.
\end{equation*}
By linearity of expectation,
\begin{equation}
    \mathcal{L} = \sum_{t=1}^T \E_{q(\mathbf{z}_{1:T})}\big[\log
    p(\mathbf{x}_t \mid \mathbf{z}_t)\big] - \sum_{t=1}^T
    \E_{q(\mathbf{z}_{1:T})}\big[r_t(\mathbf{z}_{t-1}, \mathbf{z}_t)\big].
    \label{eq:elbo_step2}
\end{equation}
Each term depends on one or two consecutive latent states only. The next
step shows how to integrate out all the others.

\subsection*{Step 3: integrating out the latent states a term does not
depend on}

\begin{proposition}\label{prop:marginalise}
For any functions $f$ and $g$,
\begin{align*}
    \E_{q(\mathbf{z}_{1:T})}\big[f(\mathbf{z}_t)\big] &=
    \E_{q(\mathbf{z}_t)}\big[f(\mathbf{z}_t)\big], \\
    \E_{q(\mathbf{z}_{1:T})}\big[g(\mathbf{z}_{t-1}, \mathbf{z}_t)\big] &=
    \E_{q(\mathbf{z}_{t-1})}\Big[\E_{q(\mathbf{z}_t \mid
    \mathbf{z}_{t-1}, \mathbf{x}_{t:T}, \mathbf{u}_{t:T})}\big[
    g(\mathbf{z}_{t-1}, \mathbf{z}_t)\big]\Big],
\end{align*}
where $q(\mathbf{z}_t)$ and $q(\mathbf{z}_{t-1})$ denote the marginals of
the inference model.
\end{proposition}

\begin{proof}
Write the expectation as an integral against
$q(\mathbf{z}_{1:T}) = \prod_{s=1}^T q(\mathbf{z}_s \mid
\mathbf{z}_{s-1}, \mathbf{x}_{s:T}, \mathbf{u}_{s:T})$.

\emph{Future states.} The integrand does not depend on $\mathbf{z}_T$
(when $t < T$), and $\mathbf{z}_T$ appears only in the last factor, which
is a probability density in $\mathbf{z}_T$. Hence
\[
    \int q(\mathbf{z}_T \mid \mathbf{z}_{T-1}, \mathbf{x}_{T},
    \mathbf{u}_{T})\,d\mathbf{z}_T = 1,
\]
and $\mathbf{z}_T$ disappears. Repeating the argument for
$\mathbf{z}_{T-1}, \dots, \mathbf{z}_{t+1}$ removes all future states.
What remains is the integral against
$q(\mathbf{z}_{1:t}) = \prod_{s=1}^{t} q(\mathbf{z}_s \mid
\mathbf{z}_{s-1}, \cdot)$, the marginal of the first $t$ states.

\emph{Past states.} Split $q(\mathbf{z}_{1:t}) =
q(\mathbf{z}_{1:t-1})\, q(\mathbf{z}_t \mid \mathbf{z}_{t-1},
\mathbf{x}_{t:T}, \mathbf{u}_{t:T})$. The inner integral over
$\mathbf{z}_t$ is the inner expectation in the statement. The resulting
function of $\mathbf{z}_{t-1}$ is then integrated against
$q(\mathbf{z}_{1:t-1})$. Integrating out $\mathbf{z}_{1:t-2}$, on which
it does not depend, gives the marginal $q(\mathbf{z}_{t-1})$ by
definition of a marginal. For $f(\mathbf{z}_t)$, integrating out
$\mathbf{z}_{1:t-1}$ in the same way gives $q(\mathbf{z}_t)$.
\end{proof}

\begin{story}[About the factor $q(\mathbf{z}_t \mid \mathbf{z}_{1:t-1},
\cdot)$]
The proof uses $q(\mathbf{z}_t \mid \mathbf{z}_{1:t-1}, \mathbf{x}_{1:T},
\mathbf{u}_{1:T}) = q(\mathbf{z}_t \mid \mathbf{z}_{t-1},
\mathbf{x}_{t:T}, \mathbf{u}_{t:T})$. This holds \emph{by construction}
of the inference model \eqref{q_phi}: the inference network is built to
receive only these inputs. D-separation is what justified that design,
because the true posterior has the same property; but in the derivation
itself, the equality is a property of $q$, not a d-separation statement.
\end{story}

\subsection*{Step 4: the result}

Apply Proposition~\ref{prop:marginalise} to \eqref{eq:elbo_step2}, with
$f = \log p(\mathbf{x}_t \mid \cdot)$ and $g = r_t$. For the second
term, the inner expectation is a KL divergence:
\begin{equation*}
    \E_{q(\mathbf{z}_t \mid \mathbf{z}_{t-1}, \mathbf{x}_{t:T},
    \mathbf{u}_{t:T})}\big[r_t(\mathbf{z}_{t-1}, \mathbf{z}_t)\big] =
    \KL\big(q(\mathbf{z}_t \mid \mathbf{z}_{t-1}, \mathbf{x}_{t:T},
    \mathbf{u}_{t:T}) \,\big\|\, p(\mathbf{z}_t \mid \mathbf{z}_{t-1},
    \mathbf{u}_{t})\big).
\end{equation*}
Therefore
\begin{equation*}
\boxed{
\begin{aligned}
    \mathcal{L} =\ & \sum_{t=1}^T \E_{q_\phi(\mathbf{z}_{t} \mid
    \mathbf{x}_{1:T}, \mathbf{u}_{1:T})} \big[\log p_\theta(\mathbf{x}_t
    \mid \mathbf{z}_{t})\big] \\
    & - \sum_{t=1}^T \E_{q_\phi(\mathbf{z}_{t-1} \mid \mathbf{x}_{1:T},
    \mathbf{u}_{1:T})} \Big[ \KL\big(q_\phi(\mathbf{z}_t \mid
    \mathbf{z}_{t-1}, \mathbf{x}_{t:T}, \mathbf{u}_{t:T}) \,\big\|\,
    p_\theta(\mathbf{z}_t \mid \mathbf{z}_{t-1}, \mathbf{u}_{t})\big)
    \Big],
\end{aligned}
}
\end{equation*}
which is \eqref{eq:dkf_elbo}. The second sum is the \emph{chain rule}
of the KL divergence: for two Markov chains,
\begin{equation*}
    \KL\big(q(\mathbf{z}_{1:T}) \,\|\, p(\mathbf{z}_{1:T})\big) =
    \sum_{t=1}^T \E_{q(\mathbf{z}_{t-1})}\Big[\KL\big(q(\mathbf{z}_t
    \mid \mathbf{z}_{t-1}) \,\|\, p(\mathbf{z}_t \mid \mathbf{z}_{t-1})
    \big)\Big].
\end{equation*}

\subsection*{Remark: iterated expectations}

Applying the proof of Proposition~\ref{prop:marginalise} to all time
steps writes the expectation under $q(\mathbf{z}_{1:T})$ as
\emph{iterated} conditional expectations,
\begin{equation*}
    \E_{q(\mathbf{z}_{1:T})}[h] =
    \E_{q(\mathbf{z}_1 \mid \mathbf{z}_0, \cdot)}\Big[
    \E_{q(\mathbf{z}_2 \mid \mathbf{z}_1, \cdot)}\big[ \cdots
    \E_{q(\mathbf{z}_T \mid \mathbf{z}_{T-1}, \cdot)}[h] \cdots \big]
    \Big].
\end{equation*}
Each inner expectation is passed to the next level as it is, with no
non-linear function applied to it, so the iterated expectations collapse
into a single expectation over the joint distribution (the tower
property). That is why ancestral sampling, one sample per level, gives
an unbiased estimate (Proposition~\ref{prop:unbiased}).

